\documentclass[11pt,a4paper]{article}
\usepackage[margin=2.5cm]{geometry}
\usepackage[T1]{fontenc}
\usepackage[utf8]{inputenc}
\usepackage[provide=*,spanish]{babel}
\usepackage{lmodern}
\usepackage{xcolor}
\usepackage{hyperref}
\usepackage{booktabs,tabularx,array}
\usepackage{longtable}
\usepackage{graphicx,float}
\usepackage{enumitem}
\hypersetup{colorlinks=true, urlcolor=blue, linkcolor=black, citecolor=black}

\begin{document}

% ===================== CARÁTULA =====================
\begin{titlepage}
\centering

{\large\bfseries UNIVERSIDAD TÉCNICA ESTATAL DE QUEVEDO\par}
\vspace{0.3cm}
{\large FACULTAD DE CIENCIAS DE LA COMPUTACIÓN\par}
\vspace{0.2cm}
{\large CARRERA DE INGENIERÍA DE SOFTWARE\par}

\vspace{2cm}
{\Large\bfseries INGENIERÍA DE REQUERIMIENTOS (ISR-401)\par}
\vspace{0.5cm}
{\large PRÁCTICA EXPERIMENTAL\par}
{\large UNIDAD IV\par}

\vspace{1cm}
{\large\itshape Inspección, gestión del cambio, trazabilidad CASE\\
y línea base del ERS FabroGym\par}

\vspace{2cm}
{\bfseries INTEGRANTES:\par}
\vspace{0.3cm}
CASTRO BAJAÑA ARIEL OMAR\\
MERA ARIAS ERICK JHAIR\\
SANCHEZ CENTENO ROSELYN ANDREINA

\vspace{1.5cm}
{\bfseries DOCENTE:\par}
\vspace{0.3cm}
PhD. Guerrero Ulloa Gleiston Cicerón

\vspace{1.5cm}
{\bfseries FECHA:\par}
\vspace{0.3cm}
05 de agosto de 2026

\vspace{2cm}
{\bfseries REPOSITORIO PÚBLICO\par}
\vspace{0.3cm}
último commit el 05/08/26 a las 23h00\\
\url{https://github.com/Emeraxs/ISR401-PFC-ERS.git}

\vfill

{\large QUEVEDO -- ECUADOR\par}
{\large 2026\par}

\end{titlepage}

% ===================== ÍNDICE =====================
\tableofcontents
\newpage

% ===================== 2. RESUMEN Y OBJETIVOS =====================
\section{Resumen y objetivos}

\subsection*{Objetivo general}

\subsection*{Objetivos específicos}
\begin{enumerate}
\item
\item
\end{enumerate}

% ===================== 3. FUNDAMENTO TEÓRICO =====================
\section{Fundamento teórico}

\subsection{Marco normativo}

\subsection{Validación de requisitos}

\subsection{Gestión de requisitos}

\subsection{Herramientas CASE para IR}

\subsection{IR en procesos ágiles}

% ===================== 4. INSPECCIÓN FAGAN =====================
\section{Inspección Fagan}

\subsection{Roles y alcance}

\subsection{Preparación individual}

\subsection{Acta de la reunión}

\subsection{Registro de defectos}

\subsection{Métricas de inspección}

% ===================== 5. CORRECCIONES APLICADAS =====================
\section{Correcciones aplicadas}

% ===================== 6. GESTIÓN DEL CAMBIO Y CCB =====================
\section{Gestión del cambio y CCB}

\subsection{RFC-01}

\subsection{RFC-02}

\subsection{RFC-03}

\subsection{Acta del CCB}

% ===================== 7. TRAZABILIDAD EN HERRAMIENTA CASE =====================
\section{Trazabilidad en herramienta CASE}

\subsection{Estructura del tablero}

\subsection{Matriz de trazabilidad}

\subsection{Evidencia}

% ===================== 8. LÍNEA BASE CON GIT =====================
\section{Línea base con Git}

% ===================== 9. COMPARATIVA DE HERRAMIENTAS CASE =====================
\section{Comparativa de herramientas CASE}

% ===================== 10. IR TRADICIONAL FRENTE A IR ÁGIL =====================
\section{IR tradicional frente a IR ágil}

% ===================== 11. CONCLUSIONES CRÍTICAS =====================
\section{Conclusiones críticas}
\begin{enumerate}
\item
\item
\item
\item
\item
\end{enumerate}

% ===================== 12. DISTRIBUCIÓN DEL TRABAJO =====================
\section{Distribución del trabajo}

% ===================== 13. REFERENCIAS =====================
\section{Referencias}

% ===================== 14. ANEXOS =====================
\section{Anexos}

\subsection*{Anexo A -- Lista de verificación de inspección}
\subsection*{Anexo B -- Registro de defectos}
\subsection*{Anexo C -- Formularios RFC y Acta del CCB}
\subsection*{Anexo D -- Exportación CSV del backlog}
\subsection*{Anexo E -- Capturas del tablero CASE y del repositorio Git}
\subsection*{Anexo F -- Referencias IEEE y declaración de uso de IA}

\end{document}
